\section{Literature Review / Theoretical Background}
Image focus detection aims to distinguish sharp regions from blurred regions
within an image, which is important for tasks such as image enhancement, object
recognition, and quality assessment. Early approaches mainly relied on
handcrafted image features, such as edge strength, gradient magnitude, or local
texture, combined with classical machine learning classifiers. These methods
are efficient and interpretable, but their performance depends heavily on the
quality of feature design and can degrade under varying image conditions.

Convolutional neural networks (CNNs) have recently become popular in image
analysis due to their ability to automatically learn hierarchical features from
raw data. CNN-based methods can be more robust than traditional approaches, but
they usually require a large number of labeled samples and substantial
computational resources. Lightweight CNN architectures reduce the number of
parameters and computational cost, but it is not clear whether they maintain an
advantage over traditional methods when training data is limited.

This study focuses on comparing traditional machine learning methods and
lightweight CNNs for segmentation-based focus detection, particularly under
low-sample conditions. The goal is to understand the trade-offs between
robustness, accuracy, and efficiency, and to identify when each approach is
most suitable.
